% Pollis Domain Packages: ArchaeologicalDetection.jl
% Compile from this folder with: julia -e 'using PollisTeX; display(compile("ArchaeologicalDetection.tex"))'
% The values of How it was checked come from the comparisons with PyTorch and Keras run on the days the text names.
% The look of a Pollis method vignette (the fonts and colours of the Pollis New Methods documents).
% Set \vignetteauthor, \vignettedate and, if not CC BY-SA 4.0, \vignettelicense after % The look of a Pollis method vignette (the fonts and colours of the Pollis New Methods documents).
% Set \vignetteauthor, \vignettedate and, if not CC BY-SA 4.0, \vignettelicense after % The look of a Pollis method vignette (the fonts and colours of the Pollis New Methods documents).
% Set \vignetteauthor, \vignettedate and, if not CC BY-SA 4.0, \vignettelicense after \input{vignette_style}.
\documentclass[10pt,a4paper]{article}

\usepackage[a4paper,left=18mm,right=18mm,top=15mm,bottom=17mm,footskip=8mm]{geometry}
\usepackage{iftex}
\ifPDFTeX\usepackage[T1]{fontenc}\fi
\usepackage[default,semibold]{sourcesanspro}
\usepackage{textcomp}
\usepackage[italic]{mathastext}
\usepackage{amsmath}
\usepackage{microtype}
\usepackage[table]{xcolor}
\usepackage{tikz}
\usepackage[most]{tcolorbox}
\usepackage{fontawesome5}
\usepackage{tabularx,booktabs}
\usepackage{enumitem}
\usepackage{fancyhdr}
\usepackage[hidelinks]{hyperref}

% ---------------------------------------------------------------- colours (as in the brochure)
\definecolor{navy}{HTML}{1C3A5E}
\definecolor{pblue}{HTML}{3B78B5}
\definecolor{amber}{HTML}{F2B33D}
\definecolor{deepamber}{HTML}{D9902A}
\definecolor{pale}{HTML}{EEF3F9}
\definecolor{ink}{HTML}{2A2F36}
\definecolor{logonavy}{HTML}{173D63}
\definecolor{logodark}{HTML}{C98A1E}
\definecolor{logoamber}{HTML}{F2B446}
\definecolor{logocream}{HTML}{FEEFD2}
\definecolor{soft}{HTML}{5E6773}
\definecolor{agree}{HTML}{2F7D4F}
\definecolor{differ}{HTML}{B23B3B}
\color{ink}

\setlength{\parindent}{0pt}
\setlength{\parskip}{4pt}
\setlist{nosep,leftmargin=*,itemsep=2pt}
\renewcommand{\arraystretch}{1.15}

\pagestyle{fancy}
\fancyhf{}
\renewcommand{\headrulewidth}{0pt}
\fancyfoot[L]{\footnotesize\color{soft}\textbf{Pollis}\enspace New Methods\enspace·\enspace\nmtheme}
\fancyfoot[R]{\footnotesize\color{soft}\thepage}
\newcommand{\nmtheme}{}
\newcommand{\vignetteauthor}{}
\newcommand{\vignettedate}{\today}
\newcommand{\vignettelicense}{\href{https://creativecommons.org/licenses/by-sa/4.0/}{CC BY-SA 4.0}}

% The Pollis logo, redrawn so it prints sharply at any size
\newcommand{\pollislogo}[1]{%
  \begin{tikzpicture}[x=#1/512,y=-#1/512,baseline=(current bounding box.center)]
    \fill[logonavy] (0,0) rectangle (512,512);
    \fill[logodark] (36,422) -- (104,256) -- (172,422) -- cycle;
    \fill[logoamber] (104,422) -- (203,180) -- (302,422) -- cycle;
    \fill[logocream] (203,422) -- (339,89) -- (476,422) -- cycle;
  \end{tikzpicture}}

% Top of the first page: logo, name; then the kicker, title and a one-line promise
\newcommand{\nmhead}[3]{%
  \renewcommand{\nmtheme}{#1}%
  \thispagestyle{fancy}%
  \begin{minipage}[c]{13mm}\pollislogo{11mm}\end{minipage}%
  \begin{minipage}[c]{\dimexpr\linewidth-13mm\relax}
    {\color{navy}\Large\bfseries Pollis}\enspace{\color{soft}New Methods}\hfill
    {\color{deepamber}\small\bfseries\MakeUppercase{#1}}
  \end{minipage}\par\vspace{7mm}
  {\color{navy}\fontsize{22}{26}\selectfont\bfseries #2}\par\vspace{4pt}
  {\color{soft}\large #3}\par\vspace{5mm}}

\newcommand{\nmsection}[1]{\par\vspace{5pt}{\color{navy}\large\bfseries #1}\par\vspace{1pt}}
\newcommand{\nmsub}[1]{\par\vspace{2pt}{\color{deepamber}\bfseries #1}\par}
\newcommand{\gt}{\,{\color{deepamber}\bfseries ›}\,}
\newcommand{\menupath}[1]{{\bfseries\color{navy}#1}}
\newcommand{\fn}[1]{\texttt{\small #1}}
% The package name in the title, in the typewriter font of the code, Latin Modern Mono
\newcommand{\pkgname}[1]{\texttt{#1}}
% Citations: \nmcite{key}{Breiman (2001)} links to the entry \nmref{key}{...} in the References
\newcommand{\nmcite}[2]{\hyperlink{ref:#1}{\color{pblue}#2}}
\newcommand{\nmref}[2]{\raisebox{1.1em}[0pt][0pt]{\hypertarget{ref:#1}{}}#2\par}
\newcommand{\nmdoi}[1]{\href{https://doi.org/#1}{\color{pblue}doi:#1}}
% Where the package is: its name and a link (its GitHub repository, once published)
\newcommand{\package}[2]{\href{#2}{\fn{#1.jl}}}

% Key facts at the top: source, missing in Julia, package, result
\newcommand{\factsheet}[4]{%
  \begin{tcolorbox}[enhanced, colback=pale, colframe=pale, boxrule=0pt, sharp corners,
    left=7pt, right=7pt, top=5pt, bottom=5pt, fontupper=\small]
  \begin{tabularx}{\linewidth}{@{}>{\leavevmode\bfseries\color{navy}}p{27mm}>{\leavevmode}X@{}}
  Tested against & #1\\
  Missing in Julia & #2\\
  Package & #3\\
  Result & #4
  \end{tabularx}
  \end{tcolorbox}}

% Boxes
\newtcolorbox{pbox}[2][]{enhanced, colback=pale, colframe=amber, boxrule=0pt,
  leftrule=3pt, sharp corners, left=7pt, right=7pt, top=6pt, bottom=6pt,
  fontupper=\small\raggedright, before upper={{\color{navy}\normalsize\bfseries #2}\par\smallskip}, #1}
\newtcolorbox{nbox}[2][]{enhanced, colback=navy, colframe=navy, boxrule=0pt,
  sharp corners, left=8pt, right=8pt, top=7pt, bottom=7pt, coltext=white,
  fontupper=\small\raggedright, before upper={{\color{amber}\normalsize\bfseries #2}\par\smallskip}, #1}

% Value tables: published beside Julia
\newenvironment{valuetable}[1]{%
  \par\small\tabularx{\linewidth}{#1}\toprule}{\bottomrule\endtabularx\par}
\newcommand{\ok}{{\color{agree}\faCheck}}
\newcommand{\no}{{\color{differ}\faTimes}}

% End of the document: credit and licence
\newcommand{\nmcredit}{%
  \par\vfill
  \begin{minipage}[b]{0.78\linewidth}\footnotesize\color{soft}
  A Pollis method, tested against the published tables on the published data.\par
  \textcopyright\ \vignetteauthor\enspace·\enspace Licensed under \vignettelicense\enspace·\enspace\vignettedate
  \end{minipage}\hfill
  \begin{minipage}[b]{0.18\linewidth}\raggedleft\pollislogo{9mm}\end{minipage}}
.
\documentclass[10pt,a4paper]{article}

\usepackage[a4paper,left=18mm,right=18mm,top=15mm,bottom=17mm,footskip=8mm]{geometry}
\usepackage{iftex}
\ifPDFTeX\usepackage[T1]{fontenc}\fi
\usepackage[default,semibold]{sourcesanspro}
\usepackage{textcomp}
\usepackage[italic]{mathastext}
\usepackage{amsmath}
\usepackage{microtype}
\usepackage[table]{xcolor}
\usepackage{tikz}
\usepackage[most]{tcolorbox}
\usepackage{fontawesome5}
\usepackage{tabularx,booktabs}
\usepackage{enumitem}
\usepackage{fancyhdr}
\usepackage[hidelinks]{hyperref}

% ---------------------------------------------------------------- colours (as in the brochure)
\definecolor{navy}{HTML}{1C3A5E}
\definecolor{pblue}{HTML}{3B78B5}
\definecolor{amber}{HTML}{F2B33D}
\definecolor{deepamber}{HTML}{D9902A}
\definecolor{pale}{HTML}{EEF3F9}
\definecolor{ink}{HTML}{2A2F36}
\definecolor{logonavy}{HTML}{173D63}
\definecolor{logodark}{HTML}{C98A1E}
\definecolor{logoamber}{HTML}{F2B446}
\definecolor{logocream}{HTML}{FEEFD2}
\definecolor{soft}{HTML}{5E6773}
\definecolor{agree}{HTML}{2F7D4F}
\definecolor{differ}{HTML}{B23B3B}
\color{ink}

\setlength{\parindent}{0pt}
\setlength{\parskip}{4pt}
\setlist{nosep,leftmargin=*,itemsep=2pt}
\renewcommand{\arraystretch}{1.15}

\pagestyle{fancy}
\fancyhf{}
\renewcommand{\headrulewidth}{0pt}
\fancyfoot[L]{\footnotesize\color{soft}\textbf{Pollis}\enspace New Methods\enspace·\enspace\nmtheme}
\fancyfoot[R]{\footnotesize\color{soft}\thepage}
\newcommand{\nmtheme}{}
\newcommand{\vignetteauthor}{}
\newcommand{\vignettedate}{\today}
\newcommand{\vignettelicense}{\href{https://creativecommons.org/licenses/by-sa/4.0/}{CC BY-SA 4.0}}

% The Pollis logo, redrawn so it prints sharply at any size
\newcommand{\pollislogo}[1]{%
  \begin{tikzpicture}[x=#1/512,y=-#1/512,baseline=(current bounding box.center)]
    \fill[logonavy] (0,0) rectangle (512,512);
    \fill[logodark] (36,422) -- (104,256) -- (172,422) -- cycle;
    \fill[logoamber] (104,422) -- (203,180) -- (302,422) -- cycle;
    \fill[logocream] (203,422) -- (339,89) -- (476,422) -- cycle;
  \end{tikzpicture}}

% Top of the first page: logo, name; then the kicker, title and a one-line promise
\newcommand{\nmhead}[3]{%
  \renewcommand{\nmtheme}{#1}%
  \thispagestyle{fancy}%
  \begin{minipage}[c]{13mm}\pollislogo{11mm}\end{minipage}%
  \begin{minipage}[c]{\dimexpr\linewidth-13mm\relax}
    {\color{navy}\Large\bfseries Pollis}\enspace{\color{soft}New Methods}\hfill
    {\color{deepamber}\small\bfseries\MakeUppercase{#1}}
  \end{minipage}\par\vspace{7mm}
  {\color{navy}\fontsize{22}{26}\selectfont\bfseries #2}\par\vspace{4pt}
  {\color{soft}\large #3}\par\vspace{5mm}}

\newcommand{\nmsection}[1]{\par\vspace{5pt}{\color{navy}\large\bfseries #1}\par\vspace{1pt}}
\newcommand{\nmsub}[1]{\par\vspace{2pt}{\color{deepamber}\bfseries #1}\par}
\newcommand{\gt}{\,{\color{deepamber}\bfseries ›}\,}
\newcommand{\menupath}[1]{{\bfseries\color{navy}#1}}
\newcommand{\fn}[1]{\texttt{\small #1}}
% The package name in the title, in the typewriter font of the code, Latin Modern Mono
\newcommand{\pkgname}[1]{\texttt{#1}}
% Citations: \nmcite{key}{Breiman (2001)} links to the entry \nmref{key}{...} in the References
\newcommand{\nmcite}[2]{\hyperlink{ref:#1}{\color{pblue}#2}}
\newcommand{\nmref}[2]{\raisebox{1.1em}[0pt][0pt]{\hypertarget{ref:#1}{}}#2\par}
\newcommand{\nmdoi}[1]{\href{https://doi.org/#1}{\color{pblue}doi:#1}}
% Where the package is: its name and a link (its GitHub repository, once published)
\newcommand{\package}[2]{\href{#2}{\fn{#1.jl}}}

% Key facts at the top: source, missing in Julia, package, result
\newcommand{\factsheet}[4]{%
  \begin{tcolorbox}[enhanced, colback=pale, colframe=pale, boxrule=0pt, sharp corners,
    left=7pt, right=7pt, top=5pt, bottom=5pt, fontupper=\small]
  \begin{tabularx}{\linewidth}{@{}>{\leavevmode\bfseries\color{navy}}p{27mm}>{\leavevmode}X@{}}
  Tested against & #1\\
  Missing in Julia & #2\\
  Package & #3\\
  Result & #4
  \end{tabularx}
  \end{tcolorbox}}

% Boxes
\newtcolorbox{pbox}[2][]{enhanced, colback=pale, colframe=amber, boxrule=0pt,
  leftrule=3pt, sharp corners, left=7pt, right=7pt, top=6pt, bottom=6pt,
  fontupper=\small\raggedright, before upper={{\color{navy}\normalsize\bfseries #2}\par\smallskip}, #1}
\newtcolorbox{nbox}[2][]{enhanced, colback=navy, colframe=navy, boxrule=0pt,
  sharp corners, left=8pt, right=8pt, top=7pt, bottom=7pt, coltext=white,
  fontupper=\small\raggedright, before upper={{\color{amber}\normalsize\bfseries #2}\par\smallskip}, #1}

% Value tables: published beside Julia
\newenvironment{valuetable}[1]{%
  \par\small\tabularx{\linewidth}{#1}\toprule}{\bottomrule\endtabularx\par}
\newcommand{\ok}{{\color{agree}\faCheck}}
\newcommand{\no}{{\color{differ}\faTimes}}

% End of the document: credit and licence
\newcommand{\nmcredit}{%
  \par\vfill
  \begin{minipage}[b]{0.78\linewidth}\footnotesize\color{soft}
  A Pollis method, tested against the published tables on the published data.\par
  \textcopyright\ \vignetteauthor\enspace·\enspace Licensed under \vignettelicense\enspace·\enspace\vignettedate
  \end{minipage}\hfill
  \begin{minipage}[b]{0.18\linewidth}\raggedleft\pollislogo{9mm}\end{minipage}}
.
\documentclass[10pt,a4paper]{article}

\usepackage[a4paper,left=18mm,right=18mm,top=15mm,bottom=17mm,footskip=8mm]{geometry}
\usepackage{iftex}
\ifPDFTeX\usepackage[T1]{fontenc}\fi
\usepackage[default,semibold]{sourcesanspro}
\usepackage{textcomp}
\usepackage[italic]{mathastext}
\usepackage{amsmath}
\usepackage{microtype}
\usepackage[table]{xcolor}
\usepackage{tikz}
\usepackage[most]{tcolorbox}
\usepackage{fontawesome5}
\usepackage{tabularx,booktabs}
\usepackage{enumitem}
\usepackage{fancyhdr}
\usepackage[hidelinks]{hyperref}

% ---------------------------------------------------------------- colours (as in the brochure)
\definecolor{navy}{HTML}{1C3A5E}
\definecolor{pblue}{HTML}{3B78B5}
\definecolor{amber}{HTML}{F2B33D}
\definecolor{deepamber}{HTML}{D9902A}
\definecolor{pale}{HTML}{EEF3F9}
\definecolor{ink}{HTML}{2A2F36}
\definecolor{logonavy}{HTML}{173D63}
\definecolor{logodark}{HTML}{C98A1E}
\definecolor{logoamber}{HTML}{F2B446}
\definecolor{logocream}{HTML}{FEEFD2}
\definecolor{soft}{HTML}{5E6773}
\definecolor{agree}{HTML}{2F7D4F}
\definecolor{differ}{HTML}{B23B3B}
\color{ink}

\setlength{\parindent}{0pt}
\setlength{\parskip}{4pt}
\setlist{nosep,leftmargin=*,itemsep=2pt}
\renewcommand{\arraystretch}{1.15}

\pagestyle{fancy}
\fancyhf{}
\renewcommand{\headrulewidth}{0pt}
\fancyfoot[L]{\footnotesize\color{soft}\textbf{Pollis}\enspace New Methods\enspace·\enspace\nmtheme}
\fancyfoot[R]{\footnotesize\color{soft}\thepage}
\newcommand{\nmtheme}{}
\newcommand{\vignetteauthor}{}
\newcommand{\vignettedate}{\today}
\newcommand{\vignettelicense}{\href{https://creativecommons.org/licenses/by-sa/4.0/}{CC BY-SA 4.0}}

% The Pollis logo, redrawn so it prints sharply at any size
\newcommand{\pollislogo}[1]{%
  \begin{tikzpicture}[x=#1/512,y=-#1/512,baseline=(current bounding box.center)]
    \fill[logonavy] (0,0) rectangle (512,512);
    \fill[logodark] (36,422) -- (104,256) -- (172,422) -- cycle;
    \fill[logoamber] (104,422) -- (203,180) -- (302,422) -- cycle;
    \fill[logocream] (203,422) -- (339,89) -- (476,422) -- cycle;
  \end{tikzpicture}}

% Top of the first page: logo, name; then the kicker, title and a one-line promise
\newcommand{\nmhead}[3]{%
  \renewcommand{\nmtheme}{#1}%
  \thispagestyle{fancy}%
  \begin{minipage}[c]{13mm}\pollislogo{11mm}\end{minipage}%
  \begin{minipage}[c]{\dimexpr\linewidth-13mm\relax}
    {\color{navy}\Large\bfseries Pollis}\enspace{\color{soft}New Methods}\hfill
    {\color{deepamber}\small\bfseries\MakeUppercase{#1}}
  \end{minipage}\par\vspace{7mm}
  {\color{navy}\fontsize{22}{26}\selectfont\bfseries #2}\par\vspace{4pt}
  {\color{soft}\large #3}\par\vspace{5mm}}

\newcommand{\nmsection}[1]{\par\vspace{5pt}{\color{navy}\large\bfseries #1}\par\vspace{1pt}}
\newcommand{\nmsub}[1]{\par\vspace{2pt}{\color{deepamber}\bfseries #1}\par}
\newcommand{\gt}{\,{\color{deepamber}\bfseries ›}\,}
\newcommand{\menupath}[1]{{\bfseries\color{navy}#1}}
\newcommand{\fn}[1]{\texttt{\small #1}}
% The package name in the title, in the typewriter font of the code, Latin Modern Mono
\newcommand{\pkgname}[1]{\texttt{#1}}
% Citations: \nmcite{key}{Breiman (2001)} links to the entry \nmref{key}{...} in the References
\newcommand{\nmcite}[2]{\hyperlink{ref:#1}{\color{pblue}#2}}
\newcommand{\nmref}[2]{\raisebox{1.1em}[0pt][0pt]{\hypertarget{ref:#1}{}}#2\par}
\newcommand{\nmdoi}[1]{\href{https://doi.org/#1}{\color{pblue}doi:#1}}
% Where the package is: its name and a link (its GitHub repository, once published)
\newcommand{\package}[2]{\href{#2}{\fn{#1.jl}}}

% Key facts at the top: source, missing in Julia, package, result
\newcommand{\factsheet}[4]{%
  \begin{tcolorbox}[enhanced, colback=pale, colframe=pale, boxrule=0pt, sharp corners,
    left=7pt, right=7pt, top=5pt, bottom=5pt, fontupper=\small]
  \begin{tabularx}{\linewidth}{@{}>{\leavevmode\bfseries\color{navy}}p{27mm}>{\leavevmode}X@{}}
  Tested against & #1\\
  Missing in Julia & #2\\
  Package & #3\\
  Result & #4
  \end{tabularx}
  \end{tcolorbox}}

% Boxes
\newtcolorbox{pbox}[2][]{enhanced, colback=pale, colframe=amber, boxrule=0pt,
  leftrule=3pt, sharp corners, left=7pt, right=7pt, top=6pt, bottom=6pt,
  fontupper=\small\raggedright, before upper={{\color{navy}\normalsize\bfseries #2}\par\smallskip}, #1}
\newtcolorbox{nbox}[2][]{enhanced, colback=navy, colframe=navy, boxrule=0pt,
  sharp corners, left=8pt, right=8pt, top=7pt, bottom=7pt, coltext=white,
  fontupper=\small\raggedright, before upper={{\color{amber}\normalsize\bfseries #2}\par\smallskip}, #1}

% Value tables: published beside Julia
\newenvironment{valuetable}[1]{%
  \par\small\tabularx{\linewidth}{#1}\toprule}{\bottomrule\endtabularx\par}
\newcommand{\ok}{{\color{agree}\faCheck}}
\newcommand{\no}{{\color{differ}\faTimes}}

% End of the document: credit and licence
\newcommand{\nmcredit}{%
  \par\vfill
  \begin{minipage}[b]{0.78\linewidth}\footnotesize\color{soft}
  A Pollis method, tested against the published tables on the published data.\par
  \textcopyright\ \vignetteauthor\enspace·\enspace Licensed under \vignettelicense\enspace·\enspace\vignettedate
  \end{minipage}\hfill
  \begin{minipage}[b]{0.18\linewidth}\raggedleft\pollislogo{9mm}\end{minipage}}


% A package Pollis needs for panels of its own, not one that reproduces a paper: its own header, key facts and credit
\renewcommand{\vignetteauthor}{Antonio Saragga Seabra}
\renewcommand{\vignettedate}{23 November 2026}
\renewcommand{\vignettelicense}{MIT}
\fancyfoot[L]{\footnotesize\color{soft}\textbf{Pollis}\enspace Domain Packages\enspace·\enspace\nmtheme}
\renewcommand{\nmhead}[3]{%
  \renewcommand{\nmtheme}{#1}%
  \thispagestyle{fancy}%
  \begin{minipage}[c]{13mm}\pollislogo{11mm}\end{minipage}%
  \begin{minipage}[c]{\dimexpr\linewidth-13mm\relax}
    {\color{navy}\Large\bfseries Pollis}\enspace{\color{soft}Domain Packages}\hfill
    {\color{deepamber}\small\bfseries\MakeUppercase{#1}}
  \end{minipage}\par\vspace{7mm}
  {\color{navy}\fontsize{22}{26}\selectfont\bfseries #2}\par\vspace{4pt}
  {\color{soft}\large #3}\par\vspace{5mm}}
\renewcommand{\factsheet}[4]{%
  \begin{tcolorbox}[enhanced, colback=pale, colframe=pale, boxrule=0pt, sharp corners,
    left=7pt, right=7pt, top=5pt, bottom=5pt, fontupper=\small]
  \begin{tabularx}{\linewidth}{@{}>{\leavevmode\bfseries\color{navy}}p{27mm}>{\leavevmode}X@{}}
  Purpose & #1\\
  Models & #2\\
  Package & #3\\
  Checked & #4
  \end{tabularx}
  \end{tcolorbox}}
\renewcommand{\nmcredit}{%
  \par\vfill
  \begin{minipage}[b]{0.78\linewidth}\footnotesize\color{soft}
  Package written and checked with Claude Code (Anthropic) in Pollis.\par
  \textcopyright\ 2026 \vignetteauthor\enspace·\enspace Licensed under \vignettelicense\enspace·\enspace\vignettedate
  \end{minipage}\hfill
  \begin{minipage}[b]{0.18\linewidth}\raggedleft\pollislogo{9mm}\par\vspace{2pt}{\footnotesize contact@pollis.org}\end{minipage}}
\newcolumntype{R}{>{\raggedleft\arraybackslash}X}
\newcolumntype{L}{>{\raggedright\arraybackslash}X}
\newcolumntype{P}[1]{>{\raggedright\arraybackslash}p{#1}}

\begin{document}
\emergencystretch=3em

\nmhead{Archaeological Prospection}
  {ArchaeologicalDetection: trained models for archaeological prospection in Julia}
  {The networks of ADAF on lidar terrain models, a random forest on Sentinel-2, and Enki on RGB imagery, without Python}

\factsheet
  {To apply trained models to a terrain model, a Sentinel-2 composite or an RGB image and draw from their output the
   features to examine as possible archaeology.}
  {The four HRNet segmentation and four Faster R-CNN object detection models of ADAF, the Automatic Detection of
   Archaeological Features tool: all archaeology, barrows, enclosures and ringforts. A random forest trained on known
   sites in a multitemporal Sentinel-2 composite. Enki, a classifier of tells on RGB imagery of 1 m, whose two dense
   layers can be trained again on known sites.}
  {\package{ArchaeologicalDetection}{https://github.com/Trumpingtons/ArchaeologicalDetection.jl}, MIT. It computes
   with NNlib.jl and DecisionTree.jl, reads ADAF's weights with Pickle.jl and Enki's with JLD2.jl and ZipArchives.jl,
   reads scenes and imagery with ArchGDAL.jl, and takes rasters read with Rasters.jl.}
  {HRNet against PyTorch on 1,048,576 cells: 22,290 reach a probability of 0.5 in Julia and 22,290 in PyTorch, and no
   cell differs on that side of 0.5; the largest difference in probability is 0.0000751. Faster R-CNN against
   torchvision: the same three boxes, their sides within 0.000062 cells and their scores within 0.000005. Enki
   against Keras on eight images: its shares within 0.00000012. Enki trained again on the protected monuments of
   one half of an area of Groningen and scored on those of the other: an AUC of 0.75 to 0.78, against 0.69 and 0.70
   for Enki as published.}

\nmsection{The models}
\nmcite{coz2026}{Čož, Corns, Curran, Kocev and Kokalj (2026)} trained neural networks to find enclosures, ringforts and
barrows in airborne laser scanning data of Ireland, labelled by experts, at cells of 0.5 m, and made them available in
ADAF. \nmcite{kostovska2025}{Kostovska, Kocev, Čož, Kokalj, Curran and Corns (2025)} published the weights: four models
for semantic segmentation, which give the probability that each cell belongs to a feature, and four for object
detection, which give boxes with scores. This package runs all eight. A trained model recognises
forms like those it was shown, so what it finds elsewhere are places to check.

\nmsection{From a terrain model to features}
\begin{valuetable}{@{}P{25mm} P{36mm} L@{}}
Step & Function & What it computes \\ \midrule
Weights & \fn{adaf\_weights} & The file of one of the eight models, 520 to 860 MB, downloaded once from Zenodo \\
Input & \fn{adaf\_input} & The simple local relief model (\nmcite{hesse2010}{Hesse, 2010}) with a radius of 10 cells, or 10 m for cells finer than 1 m, clipped to $-0.5$ to 0.5 m and stretched onto 0 to 1; 0 where it is undefined \\
Network & \fn{adaf\_model} & The HRNet or the Faster R-CNN, as the file holds, with its weights \\
Boxes & \fn{adaf\_boxes} & Faster R-CNN only: the boxes and their scores, on the tiles of the next step \\
Probability & \fn{adaf\_probability} & The network on tiles of 1024 by 1024 cells from the north-west corner, filled with 0 beyond the terrain model; for Faster R-CNN, the score of the best box over each cell \\
Features & \fn{adaf\_features} & The cells with a probability of at least the threshold, without the 8 cells along the edges of each tile (none for Faster R-CNN, with \texttt{margin = 0}), joined through their sides; each feature's area, roundness and highest probability, and the filters on area and roundness \\
\end{valuetable}

\nmsub{Use}
{\small\begin{verbatim}
using ArchaeologicalDetection, Rasters, ArchGDAL

dem = replace_missing(read(Raster("terrain.tif")), NaN32)
model = adaf_model(adaf_weights("models"; feature = :all))
probability = adaf_probability(model, dem)
found = adaf_features(probability; threshold = 0.5)
found.candidates
\end{verbatim}
}

\nmsection{The HRNet port}
ADAF builds its segmentation network with AiTLAS
(\nmcite{dimitrovski2023}{Dimitrovski, Kitanovski, Panov, Kostovska, Simidjievski and Kocev, 2023}): a High-Resolution
Network (\nmcite{wang2021}{Wang and 11 coauthors, 2021}) with 48 channels in its finest branch, as PyTorch Image Models
defines it (\nmcite{wightman2019}{Wightman, 2019}), under a head of AiTLAS's own. The port writes both with the
convolutions of NNlib.jl.
\begin{itemize}
\item \textbf{Backbone.} Two convolutions of stride 2, a layer of bottleneck blocks, then three stages of two, three
  and four branches at a quarter, an eighth, a sixteenth and a thirty-second of the resolution of the input, each
  module ending in the fusion of every branch into every other.
\item \textbf{Head.} The output of the first layer and of the four branches brought to 120 by 120 cells by nearest
  neighbour, joined into 976 channels, and projected to two outputs. The sigmoid of the second, enlarged to the tile by
  bilinear interpolation, is the probability.
\item \textbf{Weights.} Pickle.jl reads the tensors of PyTorch's checkpoint without running Python. Their axes are
  reversed for NNlib.jl.
\item \textbf{Batch normalisation.} The network is only applied, never trained, so each batch normalisation is folded
  into the weights and bias of the convolution before it.
\item \textbf{Interpolation.} Nearest and bilinear interpolation follow PyTorch's rules, in single precision and
  without aligning the corners.
\end{itemize}

\nmsection{The Faster R-CNN port}
ADAF's object detection models are torchvision's improved Faster R-CNN (\nmcite{ren2017}{Ren, He, Girshick and Sun,
2017}): a ResNet-50 (\nmcite{he2016}{He, Zhang, Ren and Sun, 2016}) under a feature pyramid
(\nmcite{lin2017}{Lin, Dollár, Girshick, He, Hariharan and Belongie, 2017}), with two classes, the feature and the
background.
\begin{itemize}
\item \textbf{Network.} The ResNet-50 and a pyramid of five levels; a region proposal network of two convolutions over
  anchors of 32 to 512 cells in three shapes; RoIAlign to 7 by 7 cells; a box head of four convolutions and a layer of
  1,024. Batch normalisations are folded as in HRNet.
\item \textbf{Image.} Each tile repeated into three bands, normalised with ImageNet's means and standard deviations
  and resized to 800 cells a side by bilinear interpolation, as torchvision does.
\item \textbf{Selection.} torchvision's: the 1,000 best anchors of each level, suppressions at overlaps of 0.7 and
  0.5, scores above 0.05, at most 100 boxes; then AiTLAS's suppression at 0.2. RoIAlign and the suppressions follow
  torchvision's code for the processor.
\item \textbf{Score map.} \fn{adaf\_probability} gives each cell the score of the best box over it, the boxes rounded
  to whole cells as ADAF rounds them, so the cells at or above a threshold are those of the boxes ADAF keeps.
\end{itemize}

\nmsection{The random forest on Sentinel-2}
\nmcite{orengo2020}{Orengo, Conesa, Garcia-Molsosa, Lobo, Green, Madella and Petrie (2020)} found hundreds of mounded
Indus settlements in Cholistan with a random forest (\nmcite{breiman2001}{Breiman, 2001}) trained in Earth Engine on
mounds already known, in a composite of the mean of thousands of Sentinel-1 and Sentinel-2 scenes. The package makes
the Sentinel-2 part of such a composite and trains the forest in Julia; unlike the networks, the model is trained
each time, on the known sites of the area.
\begin{itemize}
\item \textbf{Scenes.} \fn{sentinel\_scenes} finds the Level-2A scenes over the area in Earth Search's catalogue,
  which serves Cloud-Optimised GeoTIFFs without an account, below a share of clouds over the scene.
\item \textbf{Composite.} \fn{sentinel\_composite} warps each scene's ten bands of 10 and 20 m (\nmcite{drusch2012}{Drusch
  and 14 coauthors, 2012}) onto one grid of 10 m in the UTM zone of the area, the nearest cell for the bands of 20 m,
  reading only the cells of the area, several scenes at once. A cell is left out of a scene where its scene
  classification marks no data, saturation, cloud shadow, cloud of medium or high probability or cirrus; the
  composite is the mean of the rest, saved as a GeoTIFF.
\item \textbf{Training.} \fn{site\_cells} marks the cells of the known sites, from polygons or from points and lines
  with a radius. \fn{forest\_model} trains 128 trees of DecisionTree.jl on them and on ten times as many cells drawn
  at random at least 500 m from every site, each tree on half the cells drawn with replacement, choosing at each split
  among three of the ten bands: the defaults of Earth Engine's random forest, which the article does not say it
  changed.
\item \textbf{Probability and features.} \fn{forest\_probability} gives each cell the share of trees that vote for a
  site; \fn{forest\_features} joins the cells at or above 0.55, the article's threshold, through their sides.
\item \textbf{Differences from the article.} Level-2A, the surface reflectance, in place of the article's Level-1C;
  no Sentinel-1; negatives drawn at random in place of polygons drawn by hand; one round of training in place of
  three.
\end{itemize}

\nmsub{Use}
{\small\begin{verbatim}
scenes = sentinel_scenes(box; years = 2023:2025, cloud = 30)
composite = sentinel_composite(scenes, box; path = "composite.tif")
sites = site_cells(composite, "known-sites.gpkg"; radius = 50)
model = forest_model(composite, sites)
probability = forest_probability(model, composite)
found = forest_features(probability; threshold = 0.55)
\end{verbatim}
}

\nmsection{The Enki port}
\nmcite{rosati2025b}{Rosati (2025b)} trained Enki to tell images of cells of about 350 by 350 m that hold a tell or a
mound-like feature of the Ancient Near East from images that do not, on more than 21,500 images cut from Google and
Bing imagery (\nmcite{rosati2025a}{Rosati, 2025a}). Its network is EfficientNet-B0 (\nmcite{tan2019}{Tan and Le,
2019}) as Keras builds it, started from ImageNet, with its channels averaged and two dense layers on top.
\begin{itemize}
\item \textbf{Weights.} \fn{enki\_weights} downloads \texttt{enki0\_0.keras} from a fixed commit of Enki's repository.
  It is a ZIP file: \texttt{config.json} names the layers in order, and \texttt{model.weights.h5}, which JLD2.jl reads,
  numbers each kind of layer in that order. Their axes are reversed for NNlib.jl.
\item \textbf{Network.} Keras's scaling by ImageNet's means and standard deviations, inside the network; a stem
  convolution of stride 2; sixteen blocks of expansion, depthwise convolution, squeeze and excitation, and projection,
  with the input added back where a block keeps its size; a top convolution of 1,280 channels, averaged over the 11 by
  11 positions; a dense layer of 256 units and one of two classes. Batch normalisations, with Keras's $\epsilon$ of
  0.001, are folded into the convolutions.
\item \textbf{Padding.} The convolutions of stride 2 pad as Keras padded them for an image of 350 pixels, one cell
  more on the right and bottom than on the left and top for some, as \texttt{config.json} states.
\item \textbf{Grid.} \fn{aerial\_image} warps imagery to 1 m in the UTM zone of the area; \fn{enki\_pooled} cuts it
  into cells of 350 pixels from the north-west corner and gives each cell's 1,280 channels; \fn{enki\_probability}
  applies the dense layers, the share of the second class being the share for a tell.
\item \textbf{Training again.} \fn{enki\_finetune} trains the two dense layers from Enki's weights, the network
  beneath unchanged, as the first stage of Enki's training script does: Adam at a learning rate of 0.0001, the
  cross-entropy of the two classes, batches of 128, a dropout of 0.3, 15 epochs, a fifth of the images kept out and
  the layers of the epoch with the lowest loss on them kept. The gradients are written out by hand. Enki's second
  stage, which trains part of the network too, is left out.
\item \textbf{Features.} \fn{enki\_features} joins the cells at or above a threshold through their sides, 0.99 by
  default, the threshold of Enki's README. With \texttt{share}, it keeps instead that share of the cells with imagery,
  those of highest share first: the dense layers trained again give shares on a scale of their own, which stay below
  0.99 (see below).
\end{itemize}

\nmsub{Use}
{\small\begin{verbatim}
using ArchaeologicalDetection, Rasters, ArchGDAL

enki = enki_model(enki_weights("models"))
image = aerial_image(box; source = :esri, path = "image.tif")
cells = enki_pooled(enki, image)
probability = enki_probability(enki, cells)
found = enki_features(probability; threshold = 0.99)

points = site_points(image, "known-sites.gpkg")
examples = enki_examples(enki, image, cells, points)
tuned = enki_finetune(enki, examples.pooled, examples.labels)
probability = enki_probability(tuned.model, cells)
found = enki_features(probability; share = 0.1)
\end{verbatim}
}

\nmsection{How it was checked}
A port is faithful when it gives what the original network gives on the same input. On 22 November 2026 the HRNet
was run in Julia and in PyTorch on one tile, and on 23 November 2026 the Faster R-CNN in Julia and in torchvision on
the same tile, and Enki in Julia and in Keras on eight images. On the same day Enki was trained again on known sites of one half of
an area and scored on those of the other, to see whether training again helps.
\begin{valuetable}{@{}P{25mm} L@{}}
 & The comparison \\ \midrule
Data & The north-west 1024 by 1024 cells of a 1 m terrain model around Wandlebury Ring, south-east of Cambridge, from the LIDAR Composite DTM (\nmcite{ea2022}{Environment Agency, 2022}) \\
Input & Computed once, by \fn{adaf\_input}, and given to both networks as the same array of single-precision numbers \\
Model & All Archaeology, the same checkpoint in both \\
Reference & PyTorch 2.15.0, with the HRNet of PyTorch Image Models 1.0.31 and the forward pass and head of AiTLAS written out from its source; the checkpoint loaded with no missing and no unexpected key. For Faster R-CNN: torchvision 0.27.1 on PyTorch 2.12.1, its model loaded with the checkpoint, then AiTLAS's suppression \\
Julia & Julia 1.13.3, NNlib.jl 0.9.60 and Pickle.jl 0.3.7 (0.9.61 and 0.3.7 for Faster R-CNN), on the processor \\
\end{valuetable}

\vspace{4pt}
\begin{valuetable}{@{}L >{\raggedleft\arraybackslash}p{36mm} >{\raggedleft\arraybackslash}p{36mm}@{}}
Value & Julia & PyTorch \\ \midrule
Cells compared & 1,048,576 & 1,048,576 \\
Cells with a probability of at least 0.5 & 22,290 & 22,290 \\
Smallest probability & $8.0827025\times10^{-21}$ & $8.0827639\times10^{-21}$ \\
Largest probability & 1.0 & 1.0 \\ \midrule
Cells on different sides of 0.5 & \multicolumn{2}{r@{}}{0} \\
Largest absolute difference in probability & \multicolumn{2}{r@{}}{$7.510185\times10^{-5}$} \\
Mean absolute difference in probability & \multicolumn{2}{r@{}}{$4.380145\times10^{-8}$} \\ \midrule
Faster R-CNN: boxes found & 3 & 3 \\
Highest score & 0.721079 & 0.721084 \\
Lowest score & 0.069049 & 0.069048 \\ \midrule
Largest absolute difference in the sides of the boxes, in cells & \multicolumn{2}{r@{}}{$6.2\times10^{-5}$} \\
Largest absolute difference in score & \multicolumn{2}{r@{}}{$5\times10^{-6}$} \\
Boxes rounded to whole cells, as ADAF rounds them & \multicolumn{2}{r@{}}{the same three} \\
\end{valuetable}

\nmsub{Enki against Keras}
\begin{valuetable}{@{}P{25mm} L@{}}
 & The comparison \\ \midrule
Data & Eight images of 350 by 350 pixels of Enki's validation set (\nmcite{rosati2025a}{Rosati, 2025a}): the first four, by name, of its tells in agricultural areas, and the first four of its negatives \\
Input & The pixels read once, from 0 to 255, and given to both networks as the same array of single-precision numbers \\
Model & \texttt{enki0\_0.keras} of Enki's repository at commit e47ea9a, the same file in both \\
Reference & Keras 3.8.0, the version that saved the file, on JAX 0.12.1, the file loaded with \texttt{safe\_mode} \\
Julia & Julia 1.13.3, NNlib.jl 0.9.61, JLD2.jl 0.6.7 and ZipArchives.jl 2.6.1, on the processor \\
\end{valuetable}

\vspace{4pt}
\begin{valuetable}{@{}L >{\raggedleft\arraybackslash}p{36mm} >{\raggedleft\arraybackslash}p{36mm}@{}}
Share for a tell & Julia & Keras \\ \midrule
Tell, output\_104\_25 & 0.9517367 & 0.9517368 \\
Tell, output\_114\_33 & 0.99981886 & 0.99981886 \\
Tell, output\_118\_34 & 0.9999981 & 0.9999981 \\
Tell, output\_119\_40 & 0.9994956 & 0.9994956 \\
Negative, output\_1061\_neg & 0.00059998606 & 0.0005999906 \\
Negative, output\_1062\_neg & 0.00039327896 & 0.00039328082 \\
Negative, output\_1063\_neg & 0.0009125063 & 0.00091250805 \\
Negative, output\_1064\_neg & 0.00043141117 & 0.00043141202 \\ \midrule
Largest absolute difference in share & \multicolumn{2}{r@{}}{$1.2\times10^{-7}$} \\
Largest absolute difference in the 1,280 averaged channels & \multicolumn{2}{r@{}}{$2.6\times10^{-6}$} \\
Largest averaged channel, in absolute value & \multicolumn{2}{r@{}}{1.89} \\
\end{valuetable}

\nmsub{Enki trained again on known sites}
Enki was trained on tells of arid ground; the terps of the north of the Netherlands, mounds raised against the sea on
wet, green land, are far from that ground. The protected monuments of an area of Groningen were split at a column of
cells into a western and an eastern half, with about as many in each. \fn{enki\_finetune} was trained on the
monuments of one half, with its negatives drawn from cells of that half, and both models were scored on the
monuments of the other half, each by the cell that holds its centre; then the halves were swapped. Each training
was run with five seeds.
\begin{valuetable}{@{}P{25mm} L@{}}
 & The measurement \\ \midrule
Data & The area from 6.42 to 6.64\,°E and 53.22 to 53.35\,°N: the current aerial photograph of 25 cm of the Netherlands (\nmcite{bmn2026}{Beeldmateriaal Nederland, 2026}), warped by \fn{aerial\_image} to 15,164 by 14,967 cells of 1 m, 1,806 cells of 350 m; and the 101 archaeological monuments of the area among the protected sites of the Rijksdienst voor het Cultureel Erfgoed (\nmcite{rce2026}{2026}), in 89 cells \\
Training & \fn{enki\_examples} with its defaults: nine images of each monument of the training half, the centre and the eight points 50 m away, 423 and 480 in all; five cells for each image at least 700 m from those monuments, as many as the half holds, 508 and 648; then \fn{enki\_finetune} with its defaults \\
Scores & On the cells of the other half: the share of its monuments whose cell is among the 5, 10 or 20\% of the cells of highest share; the same at the share of cells Enki as published keeps at 0.99; and the AUC, the chance that the cell of a monument has a higher share than a cell without one \\
Julia & Julia 1.13.3, NNlib.jl 0.9.61, Rasters.jl 0.15.0 and ArchGDAL.jl 0.10.12, on the processor \\
\end{valuetable}

\vspace{4pt}
\begin{valuetable}{@{}L >{\raggedleft\arraybackslash}p{19mm} >{\raggedleft\arraybackslash}p{19mm} >{\raggedleft\arraybackslash}p{19mm} >{\raggedleft\arraybackslash}p{19mm} >{\raggedleft\arraybackslash}p{23mm}@{}}
Monuments found & 5\% & 10\% & 20\% & AUC & At 0.99's share \\ \midrule
\multicolumn{6}{@{}l}{\textit{Trained on the west, scored on the east: 54 monuments, 1,008 cells}} \\
Enki as published & 0.241 & 0.333 & 0.537 & 0.700 & 0.204 \\
Trained again, five seeds & 0.204--0.259 & 0.389--0.426 & 0.537--0.611 & 0.753--0.760 & 0.167--0.204 \\ \midrule
\multicolumn{6}{@{}l}{\textit{Trained on the east, scored on the west: 47 monuments, 798 cells}} \\
Enki as published & 0.085 & 0.170 & 0.404 & 0.692 & 0.128 \\
Trained again, five seeds & 0.255 & 0.277--0.319 & 0.489--0.532 & 0.764--0.780 & 0.277 \\
\end{valuetable}

\vspace{4pt}
In all ten trainings, the cells of the held-out monuments rank higher than with Enki as published: the AUC rises by
0.05 to 0.09. With a tenth of the cells kept, training again finds 0.39 to 0.43 of the monuments in place of 0.33 on
the east, and 0.28 to 0.32 in place of 0.17 on the west. Among the very highest cells it helps on the west only:
there, Enki as published keeps 6.5\% of the cells at 0.99, and finds 0.128 of the monuments, training again 0.277 in
as many cells; on the east, at 2.9\%, training again finds 0.167 to 0.204, against 0.204. Enki as published gives
0.99 or more to 29 cells of the eastern half and 52 of the western, Enki trained again to none and to one, which is why \fn{enki\_features} takes a share for it. The
epoch with the lowest loss on the images kept out was the fifteenth, the last, in every training.

\nmsub{What the check does not cover}
\begin{itemize}
\item \textbf{The other three models.} Barrows, enclosures and ringforts were not compared, for either network.
\item \textbf{The input image.} Both networks read the image computed in Julia, which was not compared with the image
  ADAF computes.
\item \textbf{The features.} \fn{adaf\_features} follows ADAF's rules, and its features were not compared with the
  polygons ADAF writes. For Faster R-CNN, boxes that touch across the edge of a tile become one feature, where ADAF
  keeps two.
\item \textbf{The random forest.} There is no reference to compare it with: the article's forest ran in Earth Engine,
  on other scenes and with Sentinel-1. It was run from the scenes to the features, and its composite was checked to
  lie on the grid GDAL gives the same area, but its map was not compared with the article's.
\item \textbf{Enki's grid and training again.} \fn{aerial\_image}, the grid of cells and \fn{enki\_finetune} were
  not compared with Enki's own workflow, which exports one image a cell from QGIS and trains in Keras. Training again
  was measured on one area of terps, not on tells, and the protected monuments, which carry no names, are not all
  terps.
\item \textbf{Detection itself.} Agreement with PyTorch and Keras says that the ports are faithful. It says nothing of
  how well the models detect archaeology, and it does not replicate the evaluations of ADAF's article or of Enki.
\end{itemize}

\nmsection{References}
{\footnotesize\raggedright\setlength{\parskip}{1pt}
\nmref{bmn2026}{Beeldmateriaal Nederland (2026). Luchtfoto Actueel Ortho 25 cm RGB. Dataset, served by PDOK. \href{https://service.pdok.nl/hwh/luchtfotorgb/wms/v1_0}{\color{pblue}service.pdok.nl/hwh/luchtfotorgb}}
\nmref{breiman2001}{Breiman, L. (2001). Random forests. \emph{Machine Learning}, 45(1), 5--32. \nmdoi{10.1023/A:1010933404324}}
\nmref{coz2026}{Čož, N., Corns, A., Curran, S., Kocev, D. and Kokalj, Ž. (2026). Automatic detection of archaeological features from airborne laser scanning data using deep learning: development and evaluation of the ADAF tool. \emph{Journal of Archaeological Science: Reports}, 71, 105733. \nmdoi{10.1016/j.jasrep.2026.105733}}
\nmref{dimitrovski2023}{Dimitrovski, I., Kitanovski, I., Panov, P., Kostovska, A., Simidjievski, N. and Kocev, D. (2023). AiTLAS: artificial intelligence toolbox for Earth observation. \emph{Remote Sensing}, 15(9), 2343. \nmdoi{10.3390/rs15092343}}
\nmref{drusch2012}{Drusch, M., Del Bello, U., Carlier, S., Colin, O., Fernandez, V., Gascon, F., Hoersch, B., Isola, C., Laberinti, P., Martimort, P., Meygret, A., Spoto, F., Sy, O., Marchese, F. and Bargellini, P. (2012). Sentinel-2: ESA's optical high-resolution mission for GMES operational services. \emph{Remote Sensing of Environment}, 120, 25--36. \nmdoi{10.1016/j.rse.2011.11.026}}
\nmref{ea2022}{Environment Agency (2022). LIDAR Composite Digital Terrain Model (DTM) -- 1m. Dataset. \href{https://environment.data.gov.uk/dataset/13787b9a-26a4-4775-8523-806d13af58fc}{\color{pblue}environment.data.gov.uk}}
\nmref{he2016}{He, K., Zhang, X., Ren, S. and Sun, J. (2016). Deep residual learning for image recognition. \emph{Proceedings of the IEEE Conference on Computer Vision and Pattern Recognition}, 770--778. \nmdoi{10.1109/CVPR.2016.90}}
\nmref{hesse2010}{Hesse, R. (2010). LiDAR-derived local relief models -- a new tool for archaeological prospection. \emph{Archaeological Prospection}, 17(2), 67--72. \nmdoi{10.1002/arp.374}}
\nmref{kostovska2025}{Kostovska, A., Kocev, D., Čož, N., Kokalj, Ž., Curran, S. and Corns, A. (2025). ADAF machine learning model weights. Dataset, version 1.0.0. \nmdoi{10.5281/zenodo.15848663}}
\nmref{lin2017}{Lin, T.-Y., Dollár, P., Girshick, R., He, K., Hariharan, B. and Belongie, S. (2017). Feature pyramid networks for object detection. \emph{Proceedings of the IEEE Conference on Computer Vision and Pattern Recognition}, 936--944. \nmdoi{10.1109/CVPR.2017.106}}
\nmref{orengo2020}{Orengo, H. A., Conesa, F. C., Garcia-Molsosa, A., Lobo, A., Green, A. S., Madella, M. and Petrie, C. A. (2020). Automated detection of archaeological mounds using machine-learning classification of multisensor and multitemporal satellite data. \emph{Proceedings of the National Academy of Sciences}, 117(31), 18240--18250. \nmdoi{10.1073/pnas.2005583117}}
\nmref{rce2026}{Rijksdienst voor het Cultureel Erfgoed (2026). Beschermde gebieden cultuurhistorie. Dataset, served by PDOK. \href{https://api.pdok.nl/rce/beschermde-gebieden-cultuurhistorie/ogc/v1}{\color{pblue}api.pdok.nl/rce/beschermde-gebieden-cultuurhistorie}}
\nmref{ren2017}{Ren, S., He, K., Girshick, R. and Sun, J. (2017). Faster R-CNN: towards real-time object detection with region proposal networks. \emph{IEEE Transactions on Pattern Analysis and Machine Intelligence}, 39(6), 1137--1149. \nmdoi{10.1109/TPAMI.2016.2577031}}
\nmref{rosati2025a}{Rosati, P. (2025a). Enki: AI for Archaeology -- Datasets for Automatic Site Recognition. Dataset. \nmdoi{10.5281/zenodo.14950564}}
\nmref{rosati2025b}{Rosati, P. (2025b). Enki: AI-Assistant for Archaeological Discoveries. Software and trained model. \href{https://github.com/archeorosati/enki}{\color{pblue}github.com/archeorosati/enki}}
\nmref{tan2019}{Tan, M. and Le, Q. V. (2019). EfficientNet: rethinking model scaling for convolutional neural networks. \emph{Proceedings of the 36th International Conference on Machine Learning}, PMLR 97, 6105--6114. \href{https://proceedings.mlr.press/v97/tan19a.html}{\color{pblue}proceedings.mlr.press}}
\nmref{wang2021}{Wang, J., Sun, K., Cheng, T., Jiang, B., Deng, C., Zhao, Y., Liu, D., Mu, Y., Tan, M., Wang, X., Liu, W. and Xiao, B. (2021). Deep high-resolution representation learning for visual recognition. \emph{IEEE Transactions on Pattern Analysis and Machine Intelligence}, 43(10), 3349--3364. \nmdoi{10.1109/TPAMI.2020.2983686}}
\nmref{wightman2019}{Wightman, R. (2019). PyTorch Image Models. \nmdoi{10.5281/zenodo.4414861}}
}

\nmcredit
\end{document}
